% !TeX program = xelatex
%-------------------------
% Одностраничное резюме Николая Кукина
%-------------------------

\documentclass[a4paper,10pt]{article}

% ---------- ЯЗЫК И ШРИФТЫ ----------
\usepackage{fontspec}
\usepackage{polyglossia}

\setdefaultlanguage{russian}
\setotherlanguage{english}

\IfFontExistsTF{PT Sans}
    {\setmainfont{PT Sans}}
    {\setmainfont{TeX Gyre Heros}}

% ---------- ОСНОВНЫЕ ПАКЕТЫ ----------
\usepackage[
    a4paper,
    top=1.05cm,
    bottom=1.05cm,
    left=1.15cm,
    right=1.15cm
]{geometry}

\usepackage{enumitem}
\usepackage[hidelinks]{hyperref}
\usepackage{tabularx}
\usepackage{xcolor}
\usepackage{microtype}
\usepackage{lastpage}
\usepackage{refcount}

% ---------- ЦВЕТА ----------
\definecolor{accent}{HTML}{1F4E79}
\definecolor{textgray}{HTML}{444444}

% ---------- ОБЩАЯ КОМПОНОВКА ----------
\pagestyle{empty}

\setlength{\parindent}{0pt}
\setlength{\parskip}{0pt}
\setlength{\tabcolsep}{0pt}
\setlength{\emergencystretch}{1em}

\raggedbottom
\raggedright
\urlstyle{same}

\renewcommand{\baselinestretch}{0.98}

% ---------- СПИСКИ ----------
\setlist[itemize]{
    leftmargin=15pt,
    labelsep=5pt,
    itemsep=1.2pt,
    topsep=2pt,
    parsep=0pt,
    partopsep=0pt
}

% ---------- КОМАНДЫ ----------

% Одинаковый заголовок для всех основных разделов:
% 6 pt перед заголовком;
% 2 pt между заголовком и линией;
% 3.5 pt между линией и содержимым.
\newcommand{\resumeSection}[1]{%
    \par
    \vspace{6pt}%
    {\fontsize{13}{14}\selectfont
     \bfseries
     \color{accent}
     #1\par}%
    \vspace{-8pt}%
    {\color{accent}\noindent\rule{\textwidth}{0.5pt}\par}%
    \vspace{2pt}%
}

\newcommand{\jobSeparator}{%
    \par
    \vspace{-4pt}%
    {\color{textgray!35}\noindent\rule{\textwidth}{0.4pt}\par}%
    \vspace{4pt}%
}

\newcommand{\resumeItem}[1]{%
    \item {\fontsize{9.5}{10.8}\selectfont #1}%
}

\newcommand{\resumeSubheading}[4]{%
    \begin{tabularx}{\textwidth}{@{}Xr@{}}
        {\fontsize{10.4}{11.4}\selectfont\bfseries #1}
        &
        {\fontsize{9.5}{10.5}\selectfont\bfseries #2}
        \\[-0.5pt]

        {\fontsize{9.5}{10.5}\selectfont\itshape #3}
        &
        {\fontsize{9.5}{10.5}\selectfont\itshape #4}
    \end{tabularx}
    \vspace{-2.5pt}
}

\newcommand{\resumeItemListStart}{%
    \begin{itemize}
}

\newcommand{\resumeItemListEnd}{%
    \end{itemize}
    \vspace{-2pt}
}

\newcommand{\tech}[1]{%
    \vspace{-0.5pt}
    {\fontsize{9.25}{10.3}\selectfont
    \textbf{Стек:} #1}%
}

\newcommand{\contactSeparator}{%
    \hspace{4pt}
    \textcolor{textgray}{\textbullet}
    \hspace{4pt}
}

\newcommand{\visibleLink}[2]{%
    \href{#1}{\textcolor{accent}{#2}}%
}

% ---------- ПРОВЕРКА ЧИСЛА СТРАНИЦ ----------
% После второй компиляции предупреждает, если резюме заняло больше страницы.
\AtEndDocument{%
    \ifnum\getpagerefnumber{LastPage}>1
        \PackageWarningNoLine{resume}{
            Резюме занимает больше одной страницы
        }
    \fi
}

% -------------------------------------------
%               ДОКУМЕНТ
% -------------------------------------------

\begin{document}

% ---------- ШАПКА ----------
\begin{center}
    {\fontsize{25}{27}\selectfont\bfseries Николай Кукин}

    \vspace{1pt}

    {\fontsize{13.5}{14.5}\selectfont
    \color{accent}\bfseries Backend-разработчик}

    \vspace{3pt}

    {\fontsize{10.8}{10.8}\selectfont
        \href{tel:+79916045191}{+7 991 604-51-91}
        \contactSeparator
        \href{mailto:nikolay.kukin@icloud.com}
        {nikolay.kukin@icloud.com}
        \contactSeparator
        \href{https://t.me/nikolaykukin}
        {t.me/nikolaykukin}
        \contactSeparator
        \href{https://github.com/nikolay-kukin}
        {github.com/nikolay-kukin}
    }
\end{center}

\vspace{-3pt}

% ---------- О СЕБЕ ----------
\resumeSection{О себе}

{\fontsize{9.6}{10.9}\selectfont
Backend-разработчик с коммерческим опытом создания B2B-систем и внутренних
производственных сервисов. Основной стек — Java, Spring Boot и PostgreSQL;
также разрабатываю backend на Go и сквозные сценарии с Flutter-клиентами.
Участвую в проектировании архитектуры, API, бизнес-логики и моделей данных.
Использую ИИ-инструменты для анализа и разработки.
}
\vspace{-10pt}
% ---------- ОПЫТ ----------
\resumeSection{Опыт работы}

% ---------- ЛОГИСТИЧЕСКИЙ СТАРТАП ----------
\resumeSubheading
    {Логистический стартап}
    {май 2026 — настоящее время}
    {Fullstack-разработчик}
    {Москва}

{\fontsize{9.5}{10.7}\selectfont
Внутренняя система управления производственными заявками, складскими
и логистическими процессами.
}

\resumeItemListStart
    \resumeItem{
        Разрабатываю REST API и бизнес-логику на Go для процессов отгрузки,
        возвратов, закупок и движения транспорта; реализую сквозные изменения
        в backend, мобильном приложении и desktop-dashboard.
    }

    \resumeItem{
        Реализовал идемпотентную обработку транспортных операций,
        серверное хранение статусов автомобилей и корректную доставку
        realtime-уведомлений.
    }

    \resumeItem{
        Участвовал в нормализации legacy-базы: разделении данных по схемам,
        создании справочников и устранении дублирующей логики.
    }

    \resumeItem{
        Разработал редактор BPMN-процессов с версионированием и интеграцией
        с Go API:
        \visibleLink
        {https://github.com/nikolay-kukin/redactor-bpmn}
        {github.com/nikolay-kukin/redactor-bpmn}.
    }
\resumeItemListEnd

\vspace{2pt}

\tech{Go, PostgreSQL, REST API, Docker, Flutter, WebSocket}

\jobSeparator

% ---------- ИСТОК ----------
\resumeSubheading
    {НПП «Исток» им. Шокина}
    {август 2024 — настоящее время}
    {Backend-разработчик}
    {Москва}

{\fontsize{9.5}{10.7}\selectfont
Микросервисные B2B-системы для работы с товарными каталогами, заказами,
технической документацией и инженерными проектами.
}

\resumeItemListStart
    \resumeItem{
        Спроектировал систему хранения файлов на базе MinIO и PostgreSQL:
        иерархия, версионирование, единый API и работа с техническими
        документами и изображениями.
    }

    \resumeItem{
        Разработал координатор асинхронной ИИ-обработки документов:
        разбиение на страницы, параллельные запросы к ИИ-сервису
        и поэтапное сохранение результатов.
    }

    \resumeItem{
        Реализовал MVP аналитического модуля: Kafka, ClickHouse, REST API
        и визуализация продуктовых метрик:
        \visibleLink
        {https://github.com/nikolay-kukin/analytic-module}
        {github.com/nikolay-kukin/analytic-module}.
    }

    \resumeItem{
        Разработал и нагрузочно протестировал стенд сравнения HTTP/2
        на Spring WebFlux и gRPC:
        \visibleLink
        {https://github.com/nikolay-kukin/grpc-vs-http2}
        {github.com/nikolay-kukin/grpc-vs-http2}.
    }
\resumeItemListEnd

\vspace{2pt}

\tech{
    Java 17, Spring Boot MVC/WebFlux, PostgreSQL, ClickHouse,
    MinIO, Kafka, Docker, Swagger
}

\jobSeparator

% ---------- ВЕБ-КОНСТРУКТОР ----------
\resumeSubheading
    {Индивидуальный проект «Веб-конструктор металлоконструкций»}
    {март 2024 — август 2024}
    {Fullstack-разработчик}
    {Москва}

{\fontsize{9.5}{10.7}\selectfont
Веб-приложение для расчёта конструкций и автоматического формирования
технико-коммерческих предложений:
\visibleLink
{https://github.com/nikolay-kukin/chgzabor}
{github.com/nikolay-kukin/chgzabor}.
}

\resumeItemListStart
    \resumeItem{
        Спроектировал микросервисную архитектуру и расширяемый модуль расчёта
        с добавлением новых типов изделий без изменения существующей логики.
    }

    \resumeItem{
        Разработал генератор PDF-документов на Apache PDFBox и браузерный
        интерфейс; сократил подготовку предложения с нескольких дней
        до нескольких минут.
    }
\resumeItemListEnd

\vspace{2pt}

\tech{
    Java 17, Spring Boot, Apache PDFBox, Docker,
    HTML, CSS, JavaScript
}

% ---------- ОБРАЗОВАНИЕ ----------
\resumeSection{Образование}

\begin{tabularx}{\textwidth}{@{}Xr@{}}
    {\fontsize{10}{11}\selectfont\bfseries
    РТУ МИРЭА — Информатика и вычислительная техника}
    &
    {\fontsize{9.5}{10.5}\selectfont\bfseries 2022–2026}
    \\[-0.5pt]

    \multicolumn{2}{@{}l@{}}{
        \fontsize{9.3}{10.4}\selectfont
        Бакалавриат с отличием. Автоматизация и цифровизация производств
        радиоэлектронной промышленности
    }
\end{tabularx}

\vspace{3pt}

\begin{tabularx}{\textwidth}{@{}Xr@{}}
    {\fontsize{9.3}{10.4}\selectfont
    \textbf{Профессиональная переподготовка:}
    Менеджер по информационным технологиям}
    &
    {\fontsize{9.3}{10.4}\selectfont\bfseries 2026}
    \\[0.5pt]

    {\fontsize{9.3}{10.4}\selectfont
    \textbf{Профессиональная переподготовка:}
    Технологии DevOps}
    &
    {\fontsize{9.3}{10.4}\selectfont\bfseries 2026}
\end{tabularx}

% ---------- НАВЫКИ ----------
\resumeSection{Технические навыки}

{\fontsize{9.3}{10.5}\selectfont
\textbf{Основные:}
Java 17, Spring Boot MVC/WebFlux, PostgreSQL, MinIO, Docker, Swagger

\vspace{1.5pt}

\textbf{Дополнительные:}
Go, Apache Kafka, ClickHouse, Redis, REST API, WebSocket, Apache PDFBox

\vspace{1pt}

\textbf{Базовый уровень:}
React, JavaScript, HTML, CSS, Dart, Flutter
}

% ---------- ДОПОЛНИТЕЛЬНО ----------
\resumeSection{Дополнительно}

{\fontsize{9.2}{10.4}\selectfont
Проходил дополнительное обучение по Java, тестированию, алгоритмам,
ИИ и промтингу. Имею сертификаты по профильным направлениям, а также дипломы
о высшем образовании и профессиональной переподготовке. Все подтверждающие
документы доступны по ссылке:
\visibleLink
{https://disk.yandex.ru/d/uQvg0pOUF47oiA}
{disk.yandex.ru/d/uQvg0pOUF47oiA}.
}

\end{document}